% RECOVERY-STRENGTHEN INDEPENDENT REVIEW::NESTED-AFFINE1::2026-10-08T20:58Z
\documentclass[11pt]{article}
\usepackage[margin=1in]{geometry}
\usepackage{amsmath,amssymb}
\begin{document}
\section*{Independent audit: fused affine access to one acquired evaluator}
\textbf{Global status: OPEN.} P versus NP remains open. This is an independent
algebraic and resource audit of the declared same-field restriction/orbit
model, not an identification of the complete vMS extractor or a SAT reduction.

\paragraph{Corpus reviewed.}
The task's explicit proposed construction and the original GitHub report
\texttt{2026-10-05\_0235\_LSPV-3\_ORBITBATCH1.tex},\\ blob
\texttt{766891e34f5151ed542e407d0f8bfee63322c6ba}, were read.
Repository AGENTS instructions were read. The coordinator is responsible for
whole-corpus and manuscript reconciliation; this narrow independent review
does not claim that survey. The primary author-hosted publication record for
Shaltiel--Umans (JACM 2005) was located; no uninspected extractor theorem is used.

\paragraph{PROVED/EXACT: coordinate fusion.}
Fix a field $\mathbb F_p$, dimensions $d_0=D,d_1,\ldots,d_k$, supplied
matrices $B_j\in\mathbb F_p^{d_{j-1}\times d_j}$, and supplied vectors
$c_j\in\mathbb F_p^{d_{j-1}}$. Let $E_j(z)=B_jz+c_j$. For every function
$F:\mathbb F_p^D\to\mathbb F_p$ and every $z\in\mathbb F_p^{d_k}$,
\[
 F(E_1(E_2(\cdots E_k(z))))=F(M_kz+b_k),
\]
where $M_0=I_D$, $b_0=0$, and
\[
 M_j=M_{j-1}B_j,\qquad b_j=M_{j-1}c_j+b_{j-1}.
\]
Both right-hand sides use the old $M_{j-1}$. Induction proves the identity.
No rank or invertibility hypothesis is needed for this identity. The
invertibility hypothesis in ORBITBATCH1 remains necessary for its separate
claim that a full orbit-address image equals the entire child domain.

\paragraph{One field evaluation.}
After computing the composed address, exactly one invocation of the acquired
root evaluator gives the required field symbol. For $q=\lceil\log_2p\rceil$,
a final declared Hadamard coordinate with mask $a\in\{0,1\}^q$ is
$\langle a,\operatorname{bin}_q(F(M_kz+b_k))\rangle\bmod2$.
It costs $O(q)$ further bit operations. This is a statement about a single
final coordinate. It does not suppress separately requested outputs,
repeated evaluator calls, or the acquisition of $F$.

\paragraph{PROVED/EXACT: conservative implementation bounds.}
Apply $E_k,E_{k-1},\ldots,E_1$ successively, discarding each previous
coordinate vector after its successor has been computed. This takes
\[
 O\!\left(\sum_{j=1}^k(d_{j-1}d_j+d_{j-1})\right)
\]
field operations and $O(\max_j(d_{j-1}+d_j)\log p)$ coordinate-buffer bits.
This excludes supplied descriptor storage/access and arithmetic scratch,
which must be added on the same machine. If dimensions halve at every
shrinking step and a fixed number of nonexpanding same-scale stages may
intervene, the bounds are $O(D^2)$ field operations and $O(D\log p)$
coordinate scratch. More generally the stated sums remain the correct
quantities; dimension shrinkage alone does not bound an unbounded number
of same-scale stages.

Naive matrix precomposition uses
$O(D\sum_j d_{j-1}d_j+D\sum_jd_{j-1})$ field operations, hence $O(D^3)$
under the above halving hypothesis. Retaining consecutive composition
matrices takes $O(D^2\log p)$ bits at peak; retaining only the final map
takes $O(Dd_k\log p)$ bits afterward. A subsequent query costs
$O(Dd_k+D)$ field operations. Generating $A^i$, constructing restriction
descriptors, address indexing, and retaining the path are additional
charges unless already explicitly included in a supplied representation.
For a bit-machine claim, multiply the field-operation bound by the chosen
field-arithmetic bit cost and add its workspace. Unit-cost field operations
are not unit-cost bit operations.

\paragraph{PROVED/EXACT: generic probe composition.}
If one level has nonrecursive work $w_i$ and executes at most $a_i$ calls to
its predecessor, each counted with multiplicity, then
\[
 Q_k\le Q_0\prod_{i=1}^k a_i+
 \sum_{j=1}^k w_j\prod_{i=j+1}^k a_i.
\]
Consequently, for $k=O(\log\log N)$, $a_i\le(\log N)^C$ with fixed
$C$, and uniformly $w_i=N^{o(1)}$, the bound is
$(Q_0+1)N^{o(1)}$. A root cost $Q_0=N^\gamma$ is therefore still
$N^{\gamma+o(1)}$, not $N^{o(1)}$ automatically. The hypothesis on $w_i$
must hold uniformly over all executed levels. This is a time accounting
statement; recursive ancestor state is a separate space obligation.
The affine construction improves on this generic product by composing
addresses before invoking $F$.

\paragraph{Falsification and limitations.}
Affine fusion does not apply automatically across a fresh low-degree
extension, a field change, or an intermediate Hadamard output reinterpreted
as a new polynomial. Those are different representations. Arbitrary affine
mixing preserves total degree but can increase individual-variable degrees;
thus membership in an extractor's permitted polynomial class must be
checked separately. A supplied evaluator is not a cheap-acquisition theorem.
Its setup, per-call time and live space, and the storage of all descriptors
remain charged. No identification of $p^D$ with the original SAT encoded
bit length is justified. There is no SAT return, closure, or all-input
compiler proved here, and no lower-bound endpoint is invoked.

\paragraph{COMPUTATIONAL validation.}
The adjacent deterministic script, seed 3058, checked 680 affine chains over
$p\in\{2,3,5,7\}$, including same-scale stages, with 680 actual modeled
root-evaluator invocations and all final masks for each field result.
It also checked 550 exact integer instances of the probe unrolling.
All checks passed in 0.037 seconds with a 25-second process alarm.
These checks validate implementation identities only, not SAT correctness,
resource asymptotics, or input-family coverage.

\paragraph{Duplication, adoption, and change.}
Adopt ORBITBATCH1's shared evaluator and its separation of distinct addresses
from emitted outputs. The additional recovered implementation fact is that
nested declared affine address transforms can be fused, so repeated
materialization of complete parent tables is not forced by these transforms.
This strengthens an upper-bound implementation component, while leaving
acquisition and the SAT-to-source length map OPEN. No claim in root MAIN is
silently promoted, and no retired endpoint is restored.

\paragraph{Next target.}
Supply the exact representation-changing edge of the intended recursive
construction, or prove it absent. Then charge acquisition and descriptor
generation in original SAT input length, with joint live-state accounting.

\paragraph{Publication audit of the complete coordinator report.}
The full coordinator report\\
\texttt{2026-10-08\_2102\_RECOVERY-STRENGTHEN3.tex}\\
and its main runner were subsequently read. The original NORMBLOW1 report,
blob\\ \texttt{4838e7717af7100028cba185ec2710d114f31801},\\ was fetched directly.
Its parameters $h=T^{100\eta}$ with fixed $\eta>0$, least admissible
power-of-two $r$ satisfying $h^r\ge T\operatorname{polylog}T$, and
$p\in[(hr)^{100},2(hr)^{100}]$ do imply bounded $r$ and
$\log p=O(\log T)$. Thus the coordinator's polylogarithmic coordinate
routing specialization is valid. This does not use the older report's
numerical range for the ambient exponent, whose upper endpoint would
require checking the minimum admissible value of $r$ when $\eta$ is large.

For request-dependent descriptors, the common $P$ in a $Q$-times-cost
bound must mean the maximum executed per-request routing sum, and $E,W$
must be uniform per-call evaluator bounds. Equivalently sum the actual
routing costs. The acquisition-space term must include any descriptors
that remain live during acquisition or reacquisition; scheduling initial
acquisition before descriptors is one legitimate way to avoid that
overlap. These are explicit scheduling/accounting clarifications, not
new acquisition hypotheses or a counterexample to the supplied interface.

The official CCC 2019 record independently identifies William M. Hoza
as the author of \emph{Typically-Correct Derandomization for Small Time
and Space}, DOI 10.4230/LIPIcs.CCC.2019.9, published 16 July 2019.
The actual paper's Theorem 5 and Appendix A were read: they justify
polynomial time in source length and $O(t)$ space for seed length $t$
under the stated parameters, with explicit discussion of field and
matrix construction and the final Hadamard conversion. The coordinator's
bibliographic correction and warning against treating this as cheap SAT
acquisition are supported. No identification with the declared
prime-field interface follows merely from that citation.

\bigskip\noindent\textbf{VERIFICATION SENTINEL: NESTED-AFFINE1 independent
audit complete; P versus NP OPEN.}
% END_OF_VERIFIED_REPORT::RECOVERY-STRENGTHEN3::NESTED-AFFINE1
\end{document}
